% /root, 19 November 2026. Independently audited by both other Part-A agents.
\subsection{All invariant covectors on nonordinary principally polarized surfaces}
\label{subsec:nonordinary-ppav-covectors}
\status{Proved and checked. Both the monodromy and obstruction/crystal agents
independently rederived the complete proof, checked the source scope, and
reran the retained exact local certificate without changing the root output.}

\begin{lemma}[Removing an arbitrary nilpotent length]
\label{lem:linear-covectors-any-nil-length}
\status{Proved under stated hypotheses}
Let $k$ be algebraically closed of characteristic two, let $W$ be a
finite-dimensional $k$-space with coordinates $s_1,\ldots,s_b$, and put
$P=\mathbf P(W)$ and $\mathcal B=W^*\otimes\mathcal O_P(-1)$.
Suppose a semilinear Frobenius on $\mathcal B$ has matrix
$M(s)^t$, where $M(s)=\sum_i s_iM_i$ and
$\bigcap_i\ker M_i=0$. On the open where the iterated linearized image
has maximal stable rank, let $\mathcal I$ be that unit-Frobenius image.
Every invariant $\mathbf F_2$-covector on its Lang local system is the
restriction of exactly one linear expression
\[
             \lambda_s=D(s),\qquad D\in\operatorname{End}_k(W),
             \qquad M(s)D(s)=D(s)^{[2]}.
\]
No assertion that the first Frobenius image is already stable is needed.
\end{lemma}
\begin{proof}
We first work at the generic point with field $K=k(P)$.
Choose $e$ such that the image of the $e$-fold linearized Frobenius is
the stable image $I$, and put $q=2^e$. We may increase $e$ so $e\geq1$.
An invariant covector is a $K$-linear map $f:I\to K$ satisfying
$f(Fx)=f(x)^2$. Extend it, provisionally over $K^{1/q}$, by
\[
                    \lambda(x)^q=f(F^e x),\qquad x\in\mathcal B_K.
\]
The unique $q$-th root is additive and $K$-linear in $x$.
It agrees with $f$ on $I$, and
$\lambda(Fx)=\lambda(x)^2$. Its coefficient vector raised to the
$q$-th power is a rational section
$w\in F_P^{e*}\mathcal B^*=W^{[q]}\otimes\mathcal O_P(q)$, with
\begin{equation}\label{eq:any-nil-lift-polynomial}
                    M(s)^{[q]}w=w^{[2]}.
\end{equation}
Here the brackets mean coefficientwise powers, and the equation is a
section equation in $W^{[2q]}\otimes\mathcal O_P(2q)$.
The formulas in different coefficient charts agree because a covector
on $\mathcal O(-1)$ has transition weight one.

If $w$ had a pole of largest order $h>0$ along a prime divisor, the
right side would have a coordinate of order $-2h$, whereas the left
side, whose matrix is regular, has every coordinate of order at least
$-h$. This is impossible. Normality of $P$ extends $w$ across the
remaining codimension-two set. Thus $w$ is a homogeneous polynomial
vector of degree $q$.

Write $w=\sum_{|\alpha|=q}w_\alpha s^\alpha$.
We prove inductively that every exponent occurring is divisible by
$2^j$, for $j=0,\ldots,e$. Suppose this holds for $j<e$, and choose
$\alpha$ not divisible coordinatewise by $2^{j+1}$.
It is not a pure exponent $qe_i$, so every $\alpha_i<q$.
For each $i$, the coefficient of $s_i^q s^\alpha$ on the left of
\eqref{eq:any-nil-lift-polynomial} is uniquely
$M_i^{[q]}w_\alpha$: no other pair $(i',\alpha')$ gives this monomial.
On the right, every exponent is divisible by $2^{j+1}$, so that
coefficient is zero. The common-kernel hypothesis, preserved under
Frobenius over perfect $k$, gives $w_\alpha=0$.
At $j=e$ only pure exponents remain. Taking their unique $q$-th roots
in $k$ gives $w=D(s)^{[q]}$, and then $\lambda=D(s)$ and the stated
matrix identity. Conversely that identity defines an invariant
covector. If its restriction to $I$ is zero, iteration gives
$\lambda(x)^q=\lambda(F^ex)=0$ for every $x$, so $D=0$.
\end{proof}

\begin{lemma}[The four-dimensional endomorphism algebra]
\label{lem:surface-frobenius-distributions}
\status{Proved under stated hypotheses}
Let $A/k$ be an abelian surface over an algebraically closed field of
characteristic two, and let $F=F_{A/k}:A\to A^{(2)}$.
Writing $K=\ker F$, there are canonical algebra identifications
\[
 \operatorname{End}_{A^{(2)}}(F_*\mathcal O_A)
       =\mathcal O(K)^*
       =k[\delta_1,\delta_2]/
          (\delta_i^2-\delta_i^{[2]}),
\]
where $\delta_1,\delta_2$ form a basis of invariant vector fields,
the multiplication on $\mathcal O(K)^*$ is convolution, and
$\delta_i^{[2]}$ is the invariant vector field obtained by composing
$\delta_i$ with itself. The vector-space basis is
$1,\delta_1,\delta_2,\delta_1\delta_2$.
If $f=f(A)$ is the two-rank, a basis of the invariant vector fields
can be chosen with the following relations:
\[
\begin{array}{c|cc}
 f\text{ and nilpotent type}&\delta_1^2&\delta_2^2\\\hline
 2&\delta_1&\delta_2\\
 1&\delta_1&0\\
 0,\ \text{zero map}&0&0\\
 0,\ \text{one chain}&\delta_2&0.
\end{array}
\]
\end{lemma}
\begin{proof}
The difference map gives
$A\times_{A^{(2)}}A\simeq A\times K$, since $F$ is a $K$-torsor.
Finite flat base change and adjunction therefore identify the
endomorphism space with
\[
 \operatorname{Hom}_A(F^*F_*\mathcal O_A,\mathcal O_A)
 =\operatorname{Hom}_A(\mathcal O_A\otimes_k\mathcal O(K),\mathcal O_A)
 =\mathcal O(K)^*,
\]
using $H^0(A,\mathcal O_A)=k$. Composition is convolution by the
torsor action. At the identity, smooth local parameters identify the
coordinate algebra of $K$ with $k[u,v]/(u^2,v^2)$ as an algebra.
The invariant differential operators $1,\delta_1,\delta_2,
\delta_1\delta_2$ have independent symbols dual to
$1,u,v,uv$. They therefore form a basis of its distribution algebra.
They commute, and their squares are primitive invariant derivations;
the displayed relations follow, with no additional relation by the
dimension count.

The square map on the two-dimensional space of invariant derivations
is semilinear. Its stable and nilpotent parts split over perfect $k$.
On the stable part Lang's equation supplies a basis fixed by squaring.
On a nilpotent part of dimension at most two there is either a zero
map or one chain; scaling the chain gives the displayed forms.
If the stable dimension is $a$, the distribution algebra has $2^a$
geometric points: each toral generator has values zero or one, and
nilpotent generators have value zero. But this algebra is the
coordinate algebra of the Cartier dual $K^D$.
The dual-isogeny identity gives
$K^D=\ker(V_{A^\vee/k})$; as relative Frobenius is bijective on
geometric points, this group has $2^{f(A^\vee)}=2^{f(A)}$ points.
Hence $a=f$.
The exact sources for these last identifications are
\citet[Theorem~7.5, p.~100; Proposition~7.34 and proof, p.~111;
Proposition~5.22 and proof, pp.~80--81]{EdixhovenVanDerGeerMoonenAV}.
In particular the argument concerns the whole scheme-theoretic
Frobenius kernel, not only its one geometric point.
\end{proof}

\begin{theorem}[Entire invariant quotient on every nonordinary surface]
\label{thm:nonordinary-ppav-all-covectors}
\status{Proved under stated hypotheses}
Let $(A,\Theta)$ be a principally polarized abelian surface over a
perfect field $k$ of characteristic two, with the polarization
represented by the line bundle $\mathcal O_A(\Theta)$ over $k$.
Assume $f(A)<2$ and put $L=\mathcal O_A(3\Theta)$.
For the relative Hasse--Witt local system $V_d$ on the smooth
maximal-rank open in the complete linear system $|L^d|$, for every
$d\geq1$ one has
\begin{equation}\label{eq:nonordinary-ppav-entire-covectors}
 \dim_{\mathbf F_2}(V_d^*)^{G_{\rm geom}}
 =\dim_{\mathbf F_2}(V_d^*)^{G_{\rm arith}}
 =\begin{cases}0,&d\text{ odd},\\2^{f(A)},&d\text{ even}.
 \end{cases}
\end{equation}
In particular the entire even-degree invariant quotient has dimension
one in both two-rank-zero types, and dimension two in two-rank one.
For hypersurfaces in the fixed $3\Theta$ embedding the same statements
hold whenever restriction onto $H^0(A,L^d)$ is surjective, in particular
for $d\geq\max(1,\operatorname{reg}(\mathcal I_A)-1)$.
This determines the invariant quotient, not the whole monodromy group.
\end{theorem}
\begin{proof}
We first work over $\bar k$ and choose an invariant top differential
$\omega$. Set $N=L^d$, $W=H^0(A,N)$, and $P=\mathbf P(W)$.
The full coherent bundle is
$\mathcal B=H^2(A,N^{-1})\otimes\mathcal O_P(-1)$, identified with
$W^*\otimes\mathcal O_P(-1)$ by Serre duality and $\omega$.
Its relative Euler subbundle is $\mathcal E=\Omega_P^1$.
The quotient $\mathcal B/\mathcal E$ is $H^2(A,\mathcal O_A)\otimes
\mathcal O_P$ with Frobenius zero, because $A$ is nonordinary.
Thus $\mathcal B$ and $\mathcal E$ have the same stable local system.
In the notation of the previous Cartier sections, the relative,
root-removed pairing is the $k$-bilinear map
\[
 \beta:W\otimes_kW\longrightarrow W^{(2)},\qquad
           (s,t)\longmapsto C_{\rm rel}(st\omega)/\omega^{(2)}.
\]
Its multiplication matrix $M(s)$ is linear in $s$, and the Frobenius
matrix on $\mathcal B$ is $M(s)^t$.

We check the common-kernel and descent hypotheses, including small
degrees. The vector bundle $F_*N$ is globally generated. In the
indecomposable theta case, Theorem~\ref{thm:ppav-three-theta-jets}
separates second jets for $3\Theta$; a Frobenius fibre has local
algebra $k[x,y]/(x^2,y^2)$, a quotient of the second-jet algebra.
In the product case the two degree-three elliptic bundles separate
length-two fibres, and their tensor product separates the full
Frobenius fibre. For higher $d$, multiply by a section of $L^{d-1}$
nonzero at the point: it is a unit on that finite local fibre.
The relative Cartier pairing
\[
                    F_*N\otimes F_*N\longrightarrow N^{(2)}
\]
is perfect, by the finite-Frobenius duality already checked in
\citet[Proposition~1.3.7 and Theorem~1.3.8, pp.~26--28]{BrionKumar2005}.
If $M(s)t=0$ for all $s$, fibrewise global generation and this
perfection force $t=0$. This proves the common-kernel hypothesis.

Lemma~\ref{lem:linear-covectors-any-nil-length} now represents every
invariant covector uniquely by a $D\in\operatorname{End}(W)$ such that
\begin{equation}\label{eq:nonordinary-exact-D}
                       \beta(s,D(s))=D(s)^{[2]}.
\end{equation}
Polarization in $s,t$ gives $\beta(s,Dt)=\beta(t,Ds)$.
If $s$ vanishes in one Frobenius fibre, this identity and the perfect
pairing show that $Ds$ vanishes there too. The kernel of the globally
surjective evaluation map is locally free. Thus the identity on all
closed fibres shows that $D$ descends to an
$\mathcal O_{A^{(2)}}$-linear endomorphism of $F_*N$.
No assumption on the nilpotent length was used.

If $d$ is odd, the polarization of $N$ is $[3d]\lambda_\Theta$.
Its restriction to the scheme $\ker F$ is injective, since $3d$ is odd.
\citet[Theorem~1.2(i), p.~45, with proof on pp.~43--45]{Oda1971}
therefore gives $\operatorname{End}(F_*N)=k$.
For a scalar identity the left side of
\eqref{eq:nonordinary-exact-D} is zero, since $C(\omega)=0$.
The only exact scalar lift is zero. This proves the odd-degree case.

For even $d$, write $N=F^*M$ with
$M=(L^{(2)})^{d/2}$ on $A^{(2)}$.
Projection formula gives
$\operatorname{End}(F_*N)=\operatorname{End}(F_*\mathcal O_A)$.
The operators of Lemma~\ref{lem:surface-frobenius-distributions}
act on $N$ by its Frobenius-pullback connection: in a pullback frame
they act on its coefficient function. We calculate all of them,
not merely a selected subspace.

Here are rational coordinates giving every required normal form.
All derivations annihilate the field of squares. In the zero-map
case choose $z$ with $a=\delta_1\delta_2z\ne0$; this exists because
the top distribution acts nontrivially on the generic torsor.
It is a square-field constant. Put
$x=\delta_2z/a$, $y=\delta_1z/a$. Then
$\delta_1=\partial_x$ and $\delta_2=\partial_y$.
In the one-chain case choose $z$ with $a=\delta_1^3z\ne0$ and put
$x=\delta_1^2z/a$, $y=\delta_1z/a$.
Then $\delta_1=\partial_x+x\partial_y$ and
$\delta_2=\partial_y$. In both cases $x,y$ are a two-basis and
the invariant volume can be normalized to $dx\wedge dy$:
its contraction with the invariant frame is a nonzero global constant.

In the mixed case the eigenspaces of $\delta_1^2=\delta_1$ on the
four-dimensional square-field extension have dimensions two and two;
multiplication by a nonzero weight-one element identifies them.
The nilpotent $\delta_2$ acts nontrivially on the weight-zero space.
Otherwise it acts there by zero, and on a weight-one generator $u$
as $\delta_2u=bu$; squaring gives $b^2=0$, hence $\delta_2=0$,
a contradiction. Choose weight-zero $z$ with $a=\delta_2z\ne0$.
Then $a$ is killed by both derivations, so belongs to the square field;
$y=z/a$ has weight zero and $\delta_2y=1$.
There is a nonzero weight-one $x$ killed by $\delta_2$: use any such
element if it is killed, and otherwise its nonzero $\delta_2$-image.
Then $\delta_1=x\partial_x$, $\delta_2=\partial_y$, and the volume
can be normalized to $(dx/x)\wedge dy$.

Write a general local section coefficient as
$s=a^2+xb^2+yc^2+xye^2$. On the ordered operator basis
$1,\delta_1,\delta_2,\delta_1\delta_2$, define $Q$ by
$C(sD(s)\omega)=Q(D)(s)\omega$.
Direct extraction of the odd exponents in $dx\wedge dy$ gives
\begin{equation}\label{eq:nonordinary-Cartier-table}
\begin{array}{c|cccc}
 &Q(1)&Q(\delta_1)&Q(\delta_2)&Q(\delta_1\delta_2)\\\hline
 \text{mixed}&0&0&\delta_2&\delta_1\delta_2\\
 \text{zero map}&0&0&0&\delta_1\delta_2\\
 \text{one chain}&0&\delta_2&0&\delta_1\delta_2.
\end{array}
\end{equation}
The map $Q$ is inverse-Frobenius-semilinear in scalar coefficients.
For example in the one-chain case
$\delta_1s=b^2+xc^2+(y+x^2)e^2$ and
$C(s\delta_1s\,dx\wedge dy)=(c^2+xe^2)dx\wedge dy$,
whereas $\delta_1\delta_2s=e^2$ is fixed by the same calculation.
In the mixed case the coefficient-extraction rule for $dx/x$ selects
even powers of $x$ and that for $dy$ selects odd powers of $y$;
it gives the first row. The retained exact program
\path{code/ppav_nonordinary_covector_checks.py} expands and verifies
all twelve displayed identities over the polynomial ring in
$x,y,a,b,c,e$, with the indicated Laurent differential.

Solving $D=Q(D)$ gives the two-dimensional $\mathbf F_2$-space
spanned by $\delta_2,\delta_1\delta_2$ in the mixed case, and the
one-dimensional $\mathbf F_2$-space spanned by
$\delta_1\delta_2$ in either nilpotent case.
Uniqueness in Lemma~\ref{lem:linear-covectors-any-nil-length}
shows that these are precisely the invariant covectors and that
none vanishes on the stable part.

Finally we check arithmetic descent. For two-rank zero the entire
geometric invariant space has dimension one over $\mathbf F_2$;
every Galois action on it is therefore trivial. In two-rank one the
stable line of the square map on invariant derivations has exactly
one nonzero fixed vector $\delta_1$, so it is defined over the
perfect field $k$. Its nilpotent line is likewise defined over $k$.
Choose $\omega$ over $k$ and the unique $\delta_2$ on that line with
$\omega(\delta_1,\delta_2)=1$.
Both exact operators $\delta_2$ and $\delta_1\delta_2$ are then
defined over $k$. The entire geometric invariant space is arithmetic.
The normality of the maximal-rank open ensures that passage to the
dense smooth-member open preserves the image of every finite cover.
Surjectivity of ambient restriction gives the same images by the
geometrically integral fibres of the associated linear projection;
the stated regularity bound guarantees that surjectivity.
\end{proof}

\begin{remark}[Retained exact checks]
\label{rem:nonordinary-ppav-covector-checks}
\status{Proved under stated hypotheses}
The root certificate ran in $0.09968187000049511$ seconds and is retained
as \path{outputs/root_ppav_nonordinary_covector_checks.json}.
The independent monodromy and obstruction/crystal reruns took
$0.09809181400123634$ and $0.1044188160012709$ seconds respectively.
The program verifies all sixteen local operator identities, including
the ordinary comparison row, and twelve finite illustrations of the
monomial uniqueness used in
Lemma~\ref{lem:linear-covectors-any-nil-length}.
The lemma itself is proved for every exponent and every dimension;
the finite illustrations are not substituted for that proof.
\end{remark}
